\section{Posición de la negación en la frase}
\textbf{Regla A1:} coloca \textit{niet} casi al final, justo antes del elemento que quieres negar.

\textbf{Patrones útiles}
\begin{itemize}[leftmargin=*]
  \item Oración neutra: [Sujeto] + [V2] + \ldots{} + \textit{niet}.\ \textit{Ej.: Ik lees het boek niet}.
  \item Con complementos de lugar/tiempo: pon \textit{niet} antes de ese complemento.\ \textit{Ej.: Wij gaan niet naar de les vandaag}.
  \item Con infinitivo final: \textit{niet} va antes del infinitivo.\ \textit{Ej.: Ik wil niet werken}.
  \item Con participio final: \textit{Ik heb het huis niet gezien}. (Guía inicial; se refuerza en capítulos de pasado.)
\end{itemize}

\textbf{Mini práctica}
\begin{itemize}[leftmargin=*]
  \item Inserta \textit{niet}: \textit{Ik ga morgen naar Gent}. (Niega el viaje.)
  \item Mueve \textit{niet} para negar el lugar: \textit{Zij studeert in Leiden}.
\end{itemize}
